\chapter{多重线性代数和行列式}
% \section{双线性型和二次型}
% \section{多重线性型}
% 多重线性型构成向量空间
% 多重线性型是泛函、二次型的推广
% % TODO:排列知识概览
% 代数与分析以一种截然不同的方式刻画度量这个概念
% T 的本质是换基，行列式此处的定义就是换基前后体积的比值

% 外积 <-> wedge 翻译不错

% 源于很简单的对面积的直觉-》多重线性型-》行列式

% 泛函这种用数字衡量一切的思路其实很多：群里我们用阶衡量元素，线性代数我们用秩/行列式衡量线性映射，分析里我们用范数衡量函数/向量等等

% TODO: 特征值奇异值分解的用处
\section{行列式}

% 忍俊不禁
% 交错线性形 与行列式的本质
\section{行列式的计算}
然而，批判的武器代替不了武器的批判。以下介绍行列式的一些计算方法。
\subsection{行列式大赏}
% todo: Hilbert determinant
\subsubsection{Vandermonde 行列式}
这玩意有很多种证明方法。
\href{https://en.wikipedia.org/wiki/Vandermonde_matrix#First_proof:_Polynomial_properties}{维基百科}上提供了一种基于多项式次数的方法，
很有意思。
\[
    \det(V)=\prod _{0\leq i<j\leq n}(x_{j}-x_{i})
\]
记忆这个式子也有诀窍：盯着第二行，后一项减前一项所有项的乘积就是行列式。
\subsubsection{爪形行列式}
把对角线上元素归一，分别消去首行元素，就能化简成下三角矩阵。
\begin{example}[爪形行列式]
    \[
        \textstyle
        \begin{vmatrix}
            1 & 1 & 1 & 1 \\
            1 & 2 & 0 & 0 \\
            1 & 0 & 3 & 0 \\
            1 & 0 & 0 & 4
        \end{vmatrix}=
        24
        \begin{vmatrix}
            1 & 1 & 1 & 1 \\
            \frac{1}{2} & 1 & 0 & 0 \\
            \frac{1}{3} & 0 & 1 & 0 \\
            \frac{1}{4} & 0 & 0 & 1
        \end{vmatrix}
        =24
        \begin{vmatrix}
            1-\frac{1}{2}-\frac{1}{3}-\frac{1}{4} & 0 & 0 & 0 \\
            \frac{1}{2} & 1 & 0 & 0 \\
            \frac{1}{3} & 0 & 1 & 0 \\
            \frac{1}{4} & 0 & 0 & 1
        \end{vmatrix}
        = -2
    \]
\end{example}
\subsubsection{X 型行列式}
\subsubsection{三对角行列式}
\subsubsection{差一点单位阵行列式}
分两类，一种是对角线上下相等。为\textbf{行和相等行列式}，这种全部加到一侧，然后下一行减去上一行就能计算：
\[
    \begin{vmatrix}
        b & a & \dots & a \\
        a & b & \dots & a \\
        \vdots &\vdots & \ddots & \vdots \\
        a & a & \dots & b
    \end{vmatrix}
    =
    \begin{vmatrix}
        (n-1)a+b & a & \dots & a \\
        0 & b-a & \dots & 0 \\
        \vdots & \vdots & \ddots & \vdots \\
        0 & 0 & \dots & b-a
    \end{vmatrix} \\
    = (b-a)^{n-1}\bigl((n-1)a+b\bigr).
\]
\subsubsection{循环行列式}
考虑行列式
\[
    A =
    \begin{pmatrix}
        a_1 & a_2 & a_3 & \cdots & a_n \\
        a_n & a_1 & a_2 & \cdots & a_{n-1} \\
        a_{n-1} & a_n & a_1 & \cdots & a_{n-2} \\
        \vdots & \vdots & \vdots & \ddots & \vdots \\
        a_2 & a_3 & a_4 & \cdots & a_1
    \end{pmatrix},
\]
令\(f(x) = a_1 + a_2x + a_3x^2 + \cdots + a_nx^{n-1}\)，
\(\varepsilon_0, \varepsilon_1, \dots, \varepsilon_{n-1}\) 是 \(n\) 次单位根。
设
\[
    v(\varepsilon)
    =
    \begin{pmatrix}
        1 & \varepsilon & \varepsilon^2 & \cdots & \varepsilon^{n-1}
    \end{pmatrix}^{\mathsf{T}}
\]
\[
    B =
    \begin{pmatrix}
        1 & 1 & \cdots & 1 \\
        \varepsilon_0 & \varepsilon_1 & \cdots & \varepsilon_{n-1} \\
        \varepsilon_0^2 & \varepsilon_1^2 & \cdots & \varepsilon_{n-1}^2 \\
        \vdots & \vdots & \ddots & \vdots \\
        \varepsilon_0^{n-1} & \varepsilon_1^{n-1} & \cdots
        & \varepsilon_{n-1}^{n-1}
    \end{pmatrix} = (v(\varepsilon_0), v(\varepsilon_1), \dots, v(\varepsilon_{n-1}))
\]
直接计算可得
\[
    Av(\varepsilon)
    =
    \begin{pmatrix}
        f(\varepsilon) \\
        \varepsilon f(\varepsilon) \\
        \varepsilon^2 f(\varepsilon) \\
        \vdots \\
        \varepsilon^{n-1}f(\varepsilon)
    \end{pmatrix}
    = f(\varepsilon)v(\varepsilon).
\]
因此，\(v(\varepsilon_k)\) 是 \(A\) 的特征向量，对应的特征值是
\(f(\varepsilon_k)\)。
从而
\[
    AB = B\diag\bigl(f(\varepsilon_0),f(\varepsilon_1),\dots,
    f(\varepsilon_{n-1})\bigr)
\]
\[
    \det A = \prod_{k=0}^{n-1}f(\varepsilon_k)
\]

以此纪念我大一上高等代数课的青葱岁月\UseVerb{salute}
\subsection{矩阵行列式引理}
高等代数没学好的知识，终于以线代的方式又回到了我身边\UseVerb{grin}

\begin{theorem}[矩阵行列式引理]
    \(A\) 是可逆方阵，\(\bm{u}\)、\(\bm{v}\) 是列向量，则：
    \[
        \det(A +\bm {uv}^{\mathsf{T}})=(1+\bm {v}^{\mathsf{T}}A^{-1}\bm {u} ) \det(A)
    \]
\end{theorem}

\begin{proof}
    \[
        \begin{vmatrix}
            I & \bm{u} \\
            \bm{v}^{\mathsf{T}} & 1
        \end{vmatrix} =
        \begin{vmatrix}
            I & \bm{u}\\
            0 & 1-\bm{v}^{\mathsf{T}}\bm{u}
        \end{vmatrix}
    \]
    从而\[
        \det(I +\bm {uv}^{\mathsf{T}})=(1+\bm {v}^{\mathsf{T}}\bm {u})
    \]
    取 \(\bm{u} \leftarrow A^{-1}u\)，得\[
        \det(A +\bm {uv}^{\mathsf{T}})=(1+\bm {v}^{\mathsf{T}}A^{-1}\bm {u} ) \det(A)
    \]
\end{proof}

这个定理还可以用于\(\bm{u},\bm{v}\) 是矩阵的情况。可以快速计算与\(A\) 相差不大的矩阵的行列式。
\begin{example}[差一点单位阵行列式]
    \begin{align*}
        \begin{vmatrix}
            b & a & \dots & a \\
            a & b & \dots & a \\
            \vdots &\vdots & \ddots & \vdots \\
            a & a & \dots & b
        \end{vmatrix} &= \abs{(b-a)I+a \bm{1}\bm{1}^{\mathsf{T}}} \\
        &= (b-a)^{n} \left( 1+a \bm{1}^{\mathsf{T}}\bm{1} \cdot \frac{1}{b-a} \right) \\
        &= (b-a)^{n-1} ((n-1)a+b)
    \end{align*}
\end{example}
